% !TeX root = ../../main.tex
% =====================================================================
%  Diagramas de momento de guiñada
%  Responsable:
%  Última revisión:
% =====================================================================
\chapter{Diagramas de momento de guiñada}
\label{ch:diagramas_ymd}

\begin{objetivos}
  \item Explicar qué es un diagrama de momento de guiñada y qué pregunta
        contesta que el capítulo \ref{ch:regimen_estacionario} no contesta.
  \item Construir uno: qué se barre, qué se bloquea y por qué hay que iterar.
  \item Leer en él el agarre máximo, el balance en el límite, la estabilidad y
        el control que le queda al piloto.
  \item Decidir un reparto de transferencia de carga a partir del diagrama, y
        saber qué parte de esa decisión aguanta que el modelo de neumático sea
        malo.
\end{objetivos}

\fichasesion{90 minutos}{Capítulos \ref{ch:modelado_neumatico} y
\ref{ch:regimen_estacionario}}{Ordenador con
\texttt{Codigo/dinamica/ymd.py}}

El capítulo anterior deja dos preguntas sin contestar. La primera: el gradiente
de subviraje describe el coche lejos del límite, pero ¿qué hace \emph{en} el
límite? La segunda: el cálculo de balance mira si cada eje llega o no llega,
pero no dice qué margen de maniobra le queda al piloto cuando llega.

El \textbf{diagrama de momento de guiñada} ---\emph{Yaw Moment Diagram} (YMD),
o método del momento de Milliken (MMM)--- contesta las dos a la vez
\cite{milliken_rcvd_1995}. Es, con diferencia, la herramienta de dinámica
vehicular que más información da por línea de código, y la que mejor queda en
un design event si se sabe explicar.

% ---------------------------------------------------------------------
\section{Construcción del diagrama YMD/MMM}
\label{sec:diagramas_ymd:01}
\index{diagrama de momento de guiñada}\index{YMD}\index{MMM}

\subsection{La idea: bloquear la guiñada}

En un coche real, si el momento de guiñada no es cero el coche acelera en
guiñada, cambia su velocidad de giro y el estado se mueve. Solo se pueden
sostener los estados con momento nulo, así que en régimen estacionario nunca se
ve el resto del mapa.

El truco del MMM es un experimento imaginario: \textbf{sujetar el coche en
guiñada}, como en un túnel de viento, y obligarle a mantener un ángulo de
deriva $\beta$ y un ángulo de dirección $\delta$ cualesquiera. Entonces se
pueden calcular las dos magnitudes que interesan para \emph{cualquier}
combinación, no solo para las sostenibles:

\begin{equation}
  \ay = \frac{\Fy^{\text{del}} + \Fy^{\text{tra}}}{m},
  \qquad
  N = a\,\Fy^{\text{del}} - b\,\Fy^{\text{tra}}
  \label{eq:ymd_ay_n}
\end{equation}

Representando $N$ frente a $\ay$ para toda una malla de $\beta$ y $\delta$ sale
una red de líneas: el diagrama. Los estados que el coche puede sostener de
verdad son los de $N = 0$; el resto del dibujo es lo que el coche
\emph{intentaría} hacer, y ahí está la información sobre estabilidad.

\begin{clave}
Con la guiñada bloqueada ($r = 0$) los ángulos de deriva se simplifican a
\begin{equation*}
  \slipa_{\text{del}} = \beta - \delta,
  \qquad
  \slipa_{\text{tra}} = \beta
\end{equation*}
y desaparece la velocidad de la ecuación \eqref{eq:slip_ejes}. Consecuencia
práctica: \textbf{para un coche sin aerodinámica, el diagrama es el mismo a
cualquier velocidad}. Un solo dibujo describe el coche entero.

En cuanto haya alerones esto se acaba, porque la carga aerodinámica cambia con
$V^2$: habrá que hacer un diagrama por velocidad, y ese es uno de los costes
ocultos del paquete aerodinámico que el capítulo
\ref{ch:merece_la_pena_aero} tiene que contabilizar.
\end{clave}

\subsection{El bucle que hay que iterar}

Hay una pescadilla que se muerde la cola: la transferencia de carga depende de
$\ay$, y $\ay$ depende de las fuerzas laterales, que dependen de la carga de
cada rueda. No se puede resolver de un tirón, así que se itera
(\cref{fig:ymd:bucle}). Con relajación, porque cerca del límite el punto fijo
oscila si se aplica entero.

\begin{figure}[htbp]
\centering
\begin{tikzpicture}[node distance=6mm and 9mm]
  \node[terminal] (ini) {elegir $\beta$ y $\delta$};
  \node[proceso, below=of ini] (alfa)
       {$\slipa_{\text{del}}=\beta-\delta$\\$\slipa_{\text{tra}}=\beta$};
  \node[proceso, below=of alfa] (trans)
       {transferencia de carga\\con la $\ay$ que haya};
  \node[proceso, below=of trans] (cargas)
       {carga de las cuatro ruedas};
  \node[proceso, below=of cargas] (fuerzas)
       {$\Fy$ de cada rueda\\con la Magic Formula};
  \node[proceso, below=of fuerzas] (nueva) {nueva $\ay$};
  \node[decision, below=of nueva] (conv) {¿ha convergido?};
  \node[terminal, below=of conv] (fin) {punto $(\ay,\,N)$};
  \node[nota, right=of trans] (n1)
       {la primera vuelta se entra con $\ay = 0$};
  \node[nota, right=of conv] (n2)
       {con relajación: cerca del límite, si no, oscila};

  \draw[flujo] (ini) -- (alfa);
  \draw[flujo] (alfa) -- (trans);
  \draw[flujo] (trans) -- (cargas);
  \draw[flujo] (cargas) -- (fuerzas);
  \draw[flujo] (fuerzas) -- (nueva);
  \draw[flujo] (nueva) -- (conv);
  \draw[flujo] (conv) -- node[etiq,right] {sí} (fin);
  \draw[flujo] (conv.west) -- ++(-1.5,0) |- node[etiq,left,pos=0.25] {no} (trans.west);
  \draw[aviso] (n1) -- (trans);
  \draw[aviso] (n2) -- (conv);
\end{tikzpicture}
\caption[Bucle de cálculo de un punto del diagrama]{Cómo se calcula un punto
del diagrama. El bucle existe porque la transferencia de carga depende de la
aceleración lateral y la aceleración lateral depende de la transferencia de
carga. El diagrama completo es este bucle repetido para cada combinación de
$\beta$ y $\delta$ de la malla.}
\label{fig:ymd:bucle}
\end{figure}

% ---------------------------------------------------------------------
\section{Lectura e interpretación}
\label{sec:diagramas_ymd:02}
\index{YMD!lectura}

La \cref{fig:ymd:diagrama} es el diagrama de un coche de Formula Student
genérico. Cuesta un rato entenderlo la primera vez, así que conviene leerlo por
partes.

\begin{description}[leftmargin=0pt,labelsep=0pt,style=nextline]
  \item[Las dos familias de líneas] Una familia es $\delta$ constante ---lo que
    hace el coche si el piloto no mueve el volante y el coche desliza más---; la
    otra es $\beta$ constante, es decir, cómo responde el coche al volante para
    un deslizamiento dado.
  \item[El eje horizontal: agarre] El punto más a la derecha marca la máxima
    aceleración lateral que el coche es capaz de generar. Aquí,
    \SI{1.57}{\gacc}.
  \item[La línea $N = 0$: lo sostenible] Solo los puntos donde el diagrama corta
    el eje horizontal son estados de equilibrio. El corte más alejado del origen
    es la \textbf{máxima aceleración lateral equilibrada}, que aquí vale
    \SI{1.56}{\gacc}: lo que el coche puede sostener en una curva de verdad.
  \item[Los cuatro cuadrantes] Arriba a la derecha el coche gira más de lo que
    el piloto pide (le sobra momento); abajo a la derecha, menos. Es
    sobreviraje y subviraje, pero medidos en newton-metro y no en sensaciones.
\end{description}

\begin{figure}[htbp]
\centering
\begin{tikzpicture}
\begin{axis}[cursoplot, height=7.4cm, width=0.86\linewidth,
  xlabel={$\ay$ (\si{\gacc})}, ylabel={momento de guiñada $N$ (\si{\newton\meter})},
  xmin=-1.85, xmax=1.85, ymin=-4400, ymax=4400,
  xtick={-1.5,-1,...,1.5}, ytick={-4000,-2000,0,2000,4000},
  legend pos=north west]
\addplot[fusalPrim,line width=0.7pt] table[empty line=jump]
  {Codigo/dinamica/ymd_delta.dat};
  \addlegendentry{$\delta$ constante}
\addplot[fusalSec,line width=0.7pt,densely dashed] table[empty line=jump]
  {Codigo/dinamica/ymd_beta.dat};
  \addlegendentry{$\beta$ constante}
\draw[fusalGris] (axis cs:-1.85,0) -- (axis cs:1.85,0);
\addplot[only marks,mark=*,mark size=2pt,teal!65!black]
  coordinates {(1.564,0) (-1.564,0)};
\node[cursonota,anchor=north east,text=teal!65!black]
  at (axis cs:1.82,-500) {máxima $\ay$\\equilibrada};
\node[cursonota,anchor=north east] at (axis cs:1.82,4250)
  {giro a la derecha,\\le sobra momento:\\sobrevira};
\node[cursonota,anchor=south east] at (axis cs:1.82,-4250)
  {giro a la derecha,\\le falta momento:\\subvira};
\end{axis}
\end{tikzpicture}
\caption[Diagrama de momento de guiñada]{Diagrama de momento de guiñada de un
coche de Formula Student genérico (\SI{300}{\kilogram}, reparto \num{48}/\num{52},
$\LLTD = \num{0.50}$), calculado con \texttt{ymd.py} y el neumático sintético
del capítulo \ref{ch:neumatico}. Cada línea azul es un ángulo de dirección
constante; cada línea roja discontinua, un ángulo de deriva del vehículo
constante. Solo los puntos sobre el eje horizontal son estados que el coche
puede sostener.}
\label{fig:ymd:diagrama}
\end{figure}

\begin{ojo}
El diagrama es \textbf{simétrico respecto del origen}, así que uno dado la
vuelta parece perfectamente correcto. Es la trampa de signos del apartado
\ref{sec:regimen_estacionario:01} otra vez, y aquí no se detecta a ojo: hay que
comprobarla a propósito con un caso conocido. En \texttt{ymd.py} la
comprobación es que girando el volante a la derecha salgan $\ay$ y $N$
positivos, y que un coche con el peso muy adelantado salga estable.
\end{ojo}

% ---------------------------------------------------------------------
\section{Estabilidad y control}
\label{sec:diagramas_ymd:03}
\index{estabilidad}\index{control}

Aquí está lo que no da ningún otro método. En cualquier punto del diagrama se
leen dos pendientes, y cada una responde a una pregunta distinta:

\begin{center}
\begin{tabularx}{\textwidth}{@{}lXX@{}}
\toprule
\textbf{Pendiente} & \textbf{Pregunta} & \textbf{Lo que quieres} \\
\midrule
$\pderiv{N}{\ay}$ a $\delta$ constante & Si el coche desliza más por su cuenta, ¿aparece un momento que lo endereza? & \textbf{Negativa}: estable \\
$\pderiv{N}{\delta}$ a $\beta$ constante & ¿Le queda al piloto autoridad con el volante? & \textbf{Grande}: hay control \\
\bottomrule
\end{tabularx}
\end{center}

\begin{clave}
Conviene leer la estabilidad como \textbf{pendiente en el diagrama}
($\pderiv{N}{\ay}$) y no como derivada respecto de $\beta$. El motivo es
práctico: la pendiente es lo que se ve dibujado y no cambia de signo aunque uno
se equivoque con el convenio de $\beta$, que es exactamente el error que más se
comete. Lo mismo dicho sin fórmulas: \emph{si el coche empieza a irse y las
líneas del diagrama bajan, el propio coche se recoloca; si suben, se va}.
\end{clave}

Lo importante es que \textbf{las dos cosas se miden en el límite, no en el
centro del diagrama}. Un coche puede ser perfectamente estable circulando
tranquilo y quedarse sin control en el límite, porque el eje delantero ya está
saturado y girar más el volante no añade nada. Eso, en el diagrama, se ve como
líneas de $\beta$ constante que se juntan por la derecha: el piloto mueve el
volante y el punto de trabajo apenas se mueve.

% ---------------------------------------------------------------------
\section{Uso en decisiones de diseño}
\label{sec:diagramas_ymd:04}
\index{YMD!uso en diseño}

La utilidad real del diagrama no es dibujarlo bonito una vez, sino dibujarlo
dos veces y comparar. El uso más rentable es barrer el reparto de transferencia
de carga del capítulo \ref{ch:regimen_estacionario} y mirar qué le pasa al
coche (\cref{fig:ymd:lltd}).

\begin{figure}[htbp]
\centering
\begin{tikzpicture}
\begin{axis}[cursopanel,
  xlabel={$\LLTD$}, ylabel={$\ay$ (\si{\gacc})},
  xmin=0.37, xmax=0.63, ymin=1.49, ymax=1.58,
  xtick={0.40,0.45,...,0.60}, ytick={1.50,1.52,...,1.58},
  title={(a) Cuánto agarre se usa}, legend pos=south west]
\addplot[curvaA] table[x=LLTD,y=AyTope] {Codigo/dinamica/ymd_lltd.dat};
  \addlegendentry{techo del coche}
\addplot[curvaB] table[x=LLTD,y=AyEquilibrada] {Codigo/dinamica/ymd_lltd.dat};
  \addlegendentry{equilibrada, $N=0$}
\end{axis}
\end{tikzpicture}%
\hfill
\begin{tikzpicture}
\begin{axis}[cursopanel,
  xlabel={$\LLTD$}, ylabel={$N/(m g \batalla)$ en el techo de agarre},
  xmin=0.37, xmax=0.63, ymin=-0.042, ymax=0.042,
  xtick={0.40,0.45,...,0.60}, ytick={-0.03,-0.02,...,0.03},
  title={(b) Hacia qué lado se desequilibra}]
\addplot[curvaC] table[x=LLTD,y=margen] {Codigo/dinamica/ymd_lltd.dat};
\draw[fusalGris] (axis cs:0.37,0) -- (axis cs:0.63,0);
\node[cursonota,anchor=north,text=fusalSec] at (axis cs:0.50,0.041)
     {$N>0$: le sobra momento, sobrevira};
\node[cursonota,anchor=south,text=fusalPrim] at (axis cs:0.50,-0.041)
     {$N<0$: le falta momento, subvira};
\end{axis}
\end{tikzpicture}
\caption[Efecto del reparto de transferencia de carga]{Lo que el diagrama dice
sobre el reparto de transferencia de carga. (a) El techo de agarre del coche
apenas se mueve (\SI{1.566}{\gacc} de punta a punta), pero la aceleración que
puede \emph{sostener} sí: hay un óptimo. (b) En el punto de máximo agarre queda
un momento de guiñada sin equilibrar, y su signo es el balance en el límite;
cruza cero en $\LLTD \approx \num{0.50}$. Calculado con \texttt{ymd.py}.}
\label{fig:ymd:lltd}
\end{figure}

\begin{clave}
La lectura del panel (a) es la que cambia la forma de pensar: \textbf{el reparto
de transferencia de carga casi no cambia cuánto agarre tiene el coche, sino si
el coche puede usarlo}. El techo se mueve un \SI{0.4}{\percent} entre los
extremos del barrido; la aceleración sostenible, un \SI{3.5}{\percent}. Un coche
mal repartido no es un coche sin agarre: es un coche con agarre que no puede
aprovechar porque un eje se rinde antes que el otro.
\end{clave}

Merece la pena señalar que este cruce ---$\LLTD \approx \num{0.50}$--- es el
mismo que sale del cálculo mucho más simple de la
\cref{fig:estacionario:lltd}, hecho con otro modelo y otro programa. Que dos
caminos independientes den el mismo número es la mejor validación barata que se
puede hacer, y conviene tenerla presente: \textbf{si el YMD y la hoja de
balance no coinciden, uno de los dos está mal}, y casi siempre es un signo o
una altura.

\begin{center}
\begin{tabularx}{\textwidth}{@{}Xl@{}}
\toprule
\textbf{Pregunta} & \textbf{¿Aguanta un modelo de neumático malo?} \\
\midrule
¿Qué reparto de barras equilibra el coche? & Sí \\
¿En qué dirección mover el $\LLTD$ si sobrevira? & Sí \\
¿Gano o pierdo estabilidad al bajar el $\CdG$? & Sí \\
¿Cuántas $g$ laterales da el coche? & No \\
¿Cuánto margen de control queda en el límite? & Solo cualitativamente \\
\bottomrule
\end{tabularx}
\end{center}

% ---------------------------------------------------------------------
%  Cierre del capítulo
% ---------------------------------------------------------------------

\begin{caso}[Todavía no tenemos diagrama]
Hoy no tenemos YMD del coche, y no es por falta de herramienta: el programa
está escrito y funciona. Es que no tenemos con qué alimentarlo. Faltan las dos
entradas de las que depende todo:

\begin{enumerate}
  \item El \textbf{modelo de neumático} del capítulo
        \ref{ch:modelado_neumatico}, que no existe porque nuestro Pirelli no
        está ensayado en el TTC.
  \item El \textbf{reparto de transferencia de carga}, que no se puede calcular
        sin la rigidez a balanceo del capítulo \ref{ch:regimen_estacionario}.
\end{enumerate}

Con lo que hay, el diagrama saldría, pero sería un dibujo bonito de un coche
que no es el nuestro. Lo honesto ---y lo que hay que decir en el design
event--- es que el diagrama es el paso siguiente y no un paso dado.

\medskip
Lo que sí se puede hacer ya, sin datos y sin coche, es usar
\texttt{ymd.py} con los parámetros genéricos para ver \emph{en qué dirección}
mueven las cosas: qué le pasa al balance al subir el centro de balanceo
trasero, o cuánto margen de $\LLTD$ hay antes de que el coche se vuelva
desagradable. Eso vale desde hoy y no depende de tener el neumático.

\medskip
\otrodia[Faltan el diagrama con nuestros parámetros, el $\LLTD$ elegido y la
justificación de por qué ese y no otro.]
\end{caso}

\begin{ojo}
\begin{itemize}
  \item Presentar un diagrama sin decir a qué velocidad y con qué carga
        aerodinámica está hecho. Sin aero da igual; con aero es imprescindible.
  \item Dibujar el diagrama del revés por un signo y no enterarse, porque es
        simétrico.
  \item Leer el agarre máximo en el punto más a la derecha en vez de en el corte
        con $N = 0$. El coche no puede sostener el primero.
  \item Juzgar la estabilidad en el centro del diagrama, que es donde a nadie le
        importa, en vez de en el límite.
  \item Barrer una sola variable y quedarse con el óptimo. El óptimo de $\LLTD$
        se mueve al cambiar la altura del centro de gravedad, el neumático o las
        presiones.
  \item Creerse las $g$ absolutas del diagrama. Escalan con el factor $\lamu$
        del capítulo \ref{ch:modelado_neumatico} igual que todo lo demás.
  \item Elegir el reparto que da más agarre y montarlo sin preguntar al piloto.
        El coche más rápido en el papel y el que el piloto puede llevar al
        límite no siempre son el mismo (capítulo \ref{ch:piloto_ensayos_pista}).
\end{itemize}
\end{ojo}

\begin{entregable}
\begin{enumerate}
  \item El \textbf{diagrama del coche} en una página, con los parámetros usados
        escritos al lado: masa, reparto, $\hcg$, vía, $\LLTD$ y de dónde sale el
        modelo de neumático.
  \item El \textbf{barrido de $\LLTD$}: aceleración equilibrada y balance en el
        límite frente al reparto, y el valor elegido señalado.
  \item La \textbf{comprobación cruzada} contra la hoja de balance del capítulo
        \ref{ch:regimen_estacionario}: los dos métodos deben dar el mismo
        reparto neutro.
  \item Una \textbf{tabla de sensibilidades} en el formato del capítulo
        \ref{ch:coche_como_sistema}: cuánto se mueve el balance por cada
        \SI{10}{\milli\meter} de $\hcg$, por cada punto de $\LLTD$ y por cada
        décima de bar de presión.
\end{enumerate}
\end{entregable}

\begin{juez}
\begin{enumerate}
  \item ¿Qué os dice vuestro diagrama que no os dijera ya el gradiente de
        subviraje?
  \item ¿A qué velocidad está calculado y por qué da igual (o no)?
  \item ¿Dónde está el punto de máxima aceleración lateral equilibrada?
  \item ¿El coche es estable en el límite? Enseñádmelo en el diagrama.
  \item ¿Cuánto control con el volante le queda al piloto ahí?
  \item ¿Qué $\LLTD$ habéis elegido y qué habéis descartado?
  \item ¿Coincide con lo que os dice la hoja de transferencia de carga?
\end{enumerate}
\end{juez}
